\structuralpart{ВВЕДЕНИЕ}
Полутоновое изображение представляет собой прямоугольную матрицу яркостей без цветовых каналов. Сегментацией называется разбиение его пикселей на классы. При пороговой сегментации классы определяются сравнением яркости с порогом. Такой способ применим к изображениям, для которых достаточно выделить две группы яркостей.

Цель работы --- разработать последовательный и параллельный алгоритмы сегментации полутоновых изображений, реализовать их и исследовать зависимость реального времени расчёта от числа рабочих потоков.

Для достижения цели необходимо:
\begin{itemize}
\item выбрать метод сегментации и определить входные и выходные данные;
\item разработать алгоритмы и схемы последовательного потока, диспетчера и рабочего потока;
\item реализовать алгоритмы на языке Си и выполнить функциональное тестирование;
\item измерить реальное время последовательной и параллельной реализаций при числе рабочих потоков от одного до восьмикратного числа логических ядер;
\item сравнить результаты и определить влияние размера изображения и числа рабочих потоков.
\end{itemize}
